\section{Regras derivadas para igualdade}\label{sec:regrasDerivadasParaIgualdade}

Nesta seção, extrairemos algumas regras usuais para manipular igualdade.

Os axiomas (FO16)--(FO18) afirmam, respectivamente, a reflexividade da igualdade e a compatibilidade da igualdade com símbolos funcionais e símbolos de predicado.
A partir deles, podemos justificar os raciocínios usuais ``se duas coisas são iguais, elas satisfazem as mesmas propriedades''.

Começaremos registrando uma forma conveniente dos axiomas de congruência.

\begin{metalemma}[Congruência com termos]\label{mlem:congruenciaComTermos}
    Seja $T=(\mathcal L,\Sigma)$ uma teoria de primeira ordem.

    \begin{enumerate}[label=(\roman*)]
        \item Se $f$ é um símbolo funcional $n$-ário da linguagem de $\mathcal L$, com $n>0$, e se $s_1,\dots,s_n,t_1,\dots,t_n$ são termos da linguagem de $T$, então
        \[
            T\vdash
            \left(\bigwedge_{i=1}^n s_i=t_i\right)
            \rightarrow
            f(s_1,\dots,s_n)=f(t_1,\dots,t_n).
        \]

        \item Se $P$ é um símbolo de predicado $n$-ário da linguagem de $\mathcal L $, com $n>0$, ou se $P$ é o símbolo lógico de igualdade, e se $s_1,\dots,s_n,t_1,\dots,t_n$ são termos da linguagem de $T$, então
        \[
            T\vdash
            \left(\bigwedge_{i=1}^n s_i=t_i\right)
            \rightarrow
            \bigl(P(s_1,\dots,s_n)\leftrightarrow P(t_1,\dots,t_n)\bigr).
        \]
        No caso em que $P$ é o símbolo lógico de igualdade, entende-se que $n=2$ e que $P(s_1,s_2)$ é a fórmula $s_1=s_2$.
    \end{enumerate}
\end{metalemma}

\begin{proof}
    As fórmulas acima são obtidas a partir dos axiomas (FO17) e (FO18), tomando fechos universais desses axiomas e aplicando instanciação universal aos termos indicados.
\end{proof}

Como primeira consequência, obtemos as propriedades elementares da igualdade.

\begin{metaproposition}[Propriedades elementares da igualdade]\label{mprop:propriedadesElementaresIgualdade}
    Seja $T$ uma teoria de primeira ordem e sejam $r,s,t$ termos da linguagem de $T$.
    Então:

    \begin{enumerate}
        \item $T\vdash r=r$.

        \item $T\vdash r=s\rightarrow s=r$.

        \item $T\vdash (r=s\wedge s=t)\rightarrow r=t$.
    \end{enumerate}
\end{metaproposition}

\begin{proof}
    O item (1) é obtido a partir de (FO16) por instanciação universal.

    Para o item (2), pelo item (2) do Metalema~\ref{mlem:congruenciaComTermos}, aplicado ao símbolo lógico de igualdade, temos
    \[
        T\vdash
        (r=s\wedge r=r)\rightarrow(r=r\leftrightarrow s=r).
    \]
    Como $T\vdash r=r$, segue, por raciocínio tautológico, que
    \[
        T\vdash r=s\rightarrow s=r.
    \]

    Para o item (3), pelo item (2), temos
    \[
        T\vdash r=s\rightarrow s=r.
    \]
    Logo, por raciocínio  tautológico,
    \[
        T\vdash (r=s\wedge s=t)\rightarrow(s=r\wedge s=t).
    \]
    Novamente pelo item (2) do Metalema~\ref{mlem:congruenciaComTermos}, aplicado ao símbolo lógico de igualdade,
    \[
        T\vdash
        (s=r\wedge s=t)\rightarrow(s=s\leftrightarrow r=t).
    \]
    Como $T\vdash s=s$, segue que
    \[
        T\vdash (s=r\wedge s=t)\rightarrow r=t.
    \]
    Portanto,
    \[
        T\vdash (r=s\wedge s=t)\rightarrow r=t.
    \]
\end{proof}

Em particular, podemos usar simetria e transitividade da igualdade como regras derivadas:
\[
    \frac{T\vdash r=s}{T\vdash s=r}
    \qquad\text{e}\qquad
    \frac{T\vdash r=s \quad T\vdash s=t}{T\vdash r=t}.
\]

Passamos agora à substituição de termos iguais dentro de termos.

\begin{metalemma}[Substituição de iguais em termos]\label{mlem:substituicaoIguaisTermos}
    Seja $T$ uma teoria de primeira ordem.
    Seja $x$ uma variável, sejam $s$ e $t$ termos da linguagem de $T$, e seja $r$ um termo da linguagem de $T$.

    Então
    \[
        T\vdash s=t\rightarrow r[x/s]=r[x/t].
    \]
\end{metalemma}

\begin{proof}
    Procedemos por indução sobre a estrutura do termo $r$.

    Se $r$ é a variável $x$, então $r[x/s]$ é $s$ e $r[x/t]$ é $t$.
    Logo, a conclusão é
    \[
        T\vdash s=t\rightarrow s=t,
    \]
    que é uma instância de tautologia.

    Se $r$ é uma variável diferente de $x$, ou se $r$ é uma constante, então $r[x/s]$ e $r[x/t]$ são ambos $r$.
    Como $T\vdash r=r$, segue que
    \[
        T\vdash s=t\rightarrow r=r.
    \]

    Finalmente, suponha que $r$ é $f(r_1,\dots,r_n)$, em que $f$ é um símbolo funcional $n$-ário e $n>0$.
    Pela hipótese de indução, para cada $i\in\{1,\dots,n\}$,
    \[
        T\vdash s=t\rightarrow r_i[x/s]=r_i[x/t].
    \]
    Por raciocínio proposicional, obtemos
    \[
        T\vdash
        s=t\rightarrow
        \bigwedge_{i=1}^n r_i[x/s]=r_i[x/t].
    \]
    Pelo Metalema~\ref{mlem:congruenciaComTermos},
    \[
        T\vdash
        \left(\bigwedge_{i=1}^n r_i[x/s]=r_i[x/t]\right)
        \rightarrow
        f(r_1[x/s],\dots,r_n[x/s])
        =
        f(r_1[x/t],\dots,r_n[x/t]).
    \]
    Portanto,
    \[
        T\vdash
        s=t\rightarrow r[x/s]=r[x/t].
    \]
\end{proof}

Agora veremos que a mesma ideia vale para fórmulas atômicas.

\begin{metalemma}[Substituição de iguais em fórmulas atômicas]\label{mlem:substituicaoIguaisAtomicas}
    Seja $T$ uma teoria de primeira ordem.
    Seja $x$ uma variável, sejam $s$ e $t$ termos da linguagem de $T$, e seja $\alpha$ uma fórmula atômica da linguagem de $T$.

    Então
    \[
        T\vdash s=t\rightarrow(\alpha[x/s]\leftrightarrow\alpha[x/t]).
    \]
\end{metalemma}

\begin{proof}
    Se $\alpha$ é um símbolo de predicado $0$-ário, então $\alpha[x/s]$ e $\alpha[x/t]$ são ambas $\alpha$.
    Logo,
    \[
        T\vdash s=t\rightarrow(\alpha\leftrightarrow\alpha).
    \]

    Suponha que $\alpha$ é $P(r_1,\dots,r_n)$, em que $P$ é um símbolo de predicado $n$-ário e $n>0$.
    Pelo Metalema~\ref{mlem:substituicaoIguaisTermos}, para cada $i\in\{1,\dots,n\}$,
    \[
        T\vdash s=t\rightarrow r_i[x/s]=r_i[x/t].
    \]
    Assim, por raciocínio proposicional,
    \[
        T\vdash
        s=t\rightarrow
        \bigwedge_{i=1}^n r_i[x/s]=r_i[x/t].
    \]
    Pelo Metalema~\ref{mlem:congruenciaComTermos},
    \[
        T\vdash
        \left(\bigwedge_{i=1}^n r_i[x/s]=r_i[x/t]\right)
        \rightarrow
        \bigl(P(r_1[x/s],\dots,r_n[x/s])
        \leftrightarrow
        P(r_1[x/t],\dots,r_n[x/t])\bigr).
    \]
    Portanto,
    \[
        T\vdash s=t\rightarrow(\alpha[x/s]\leftrightarrow\alpha[x/t]).
    \]

    Finalmente, se $\alpha$ é uma fórmula equacional $r_1=r_2$, procedemos da mesma forma, usando o item (2) do Metalema~\ref{mlem:congruenciaComTermos} aplicado ao símbolo lógico de igualdade.
\end{proof}

Precisaremos do seguinte.

\begin{metalemma}[Quantificação de equivalências sob hipótese]\label{mlem:quantificacaoEquivalenciasSobHipotese}
    Seja $T$ uma teoria de primeira ordem.
    Seja $y$ uma variável e sejam $\theta,\alpha,\beta$ fórmulas da linguagem de $T$.
    Suponha que $y$ não ocorra livre em $\theta$.

    Então, para $Q$ sendo $\forall$ ou $\exists$, vale a seguinte regra derivada:
    \[
        \frac{T\vdash \theta\rightarrow(\alpha\leftrightarrow\beta)}
             {T\vdash \theta\rightarrow(Qy\,\alpha\leftrightarrow Qy\,\beta)}.
    \]
\end{metalemma}

\begin{proof}
    Para $\forall$, observe primeiro que $T\vdash \theta\rightarrow\forall y \theta$, pela introdução do quantificador universal.
    Agora note que, da generalização da hipótese do metalema, $T\vdash \forall y(\theta\rightarrow (\alpha\leftrightarrow\beta))$, e que, portanto, $T\vdash \forall y \theta\rightarrow \forall y(\alpha\leftrightarrow\beta)$.
    Assim, pela transitividade de $\rightarrow$, $T\vdash \theta\rightarrow \forall y(\alpha\leftrightarrow\beta)$.

    Como a passagem de $\forall y(\alpha\leftrightarrow\beta)$ para $\forall y\,\alpha\leftrightarrow\forall y\,\beta$ é uma regra derivada das regras para quantificadores, segue, pela transitividade de $\rightarrow$, que
    \[
        T\vdash \theta\rightarrow(\forall y\,\alpha\leftrightarrow\forall y\,\beta).'
    \]
    
    Para $\exists$, temos da hipótese que $T\vdash \theta\rightarrow(\neg \alpha\leftrightarrow\neg \beta)$ por raciocínio tautológico.
    Da primeira parte, temos $T\vdash \theta\rightarrow(\forall y\,\neg \alpha\leftrightarrow \forall y\,\neg \beta)$.
    Novamente de raciocínio tautológico, $T\vdash \theta\rightarrow(\neg\forall y\,\neg \alpha\leftrightarrow \neg\forall y\,\neg \beta)$.
    
    Como $T\vdash \neg\forall y\,\neg \alpha\leftrightarrow \exists y\,\alpha$ e $T\vdash \neg\forall y\,\neg \beta\leftrightarrow \exists y\,\beta$, segue que $T\vdash \theta\rightarrow(\exists y\,\alpha\leftrightarrow \exists y\,\beta)$.
\end{proof}

Agora estamos prontos para demonstrar a Lei de Leibniz.

\begin{metatheorem}[Lei da igualdade de Leibniz]\label{mthm:leiIgualdadeLeibniz}
    Seja $T$ uma teoria de primeira ordem.
    Seja $x$ uma variável, sejam $s$ e $t$ termos da linguagem de $T$, e seja $\varphi$ uma fórmula da linguagem de $T$.

    Suponha que $x$ seja substituível por $s$ e por $t$ em $\varphi$.
    Então
    \[
        T\vdash s=t\rightarrow(\varphi[x/s]\leftrightarrow\varphi[x/t]).
    \]
\end{metatheorem}

\begin{proof}
    Procedemos por indução sobre a estrutura de $\varphi$.

    Se $\varphi$ é atômica, a conclusão é exatamente o Metalema~\ref{mlem:substituicaoIguaisAtomicas}.

    Se $\varphi$ é $\neg\psi$, então, por hipótese de indução,
    \[
        T\vdash s=t\rightarrow(\psi[x/s]\leftrightarrow\psi[x/t]).
    \]
    
    Ademais, $T\vdash (\psi[x/s]\leftrightarrow\psi[x/t])\rightarrow (\neg\psi[x/s]\leftrightarrow\neg\psi[x/t])$.
    Pela transitividade de $\rightarrow$, obtemos
    \[
        T\vdash s=t\rightarrow(\neg\psi[x/s]\leftrightarrow\neg\psi[x/t]).
    \]

    Se $\varphi$ é $\psi\square\chi$, em que $\square$ é um conectivo binário, então, pelas hipóteses de indução aplicadas a $\psi$ e a $\chi$,
    \[
        T\vdash s=t\rightarrow(\psi[x/s]\leftrightarrow\psi[x/t])
    \]
    e
    \[
        T\vdash s=t\rightarrow(\chi[x/s]\leftrightarrow\chi[x/t]).
    \]
    Ademais, $T\vdash
    ((\psi[x/s]\leftrightarrow\psi[x/t])\wedge(\chi[x/s]\leftrightarrow\chi[x/t]))\rightarrow ((\psi[x/s]\square\chi[x/s])\leftrightarrow(\psi[x/t]\square\chi[x/t]))$.
    Assim, 
    \[
        T\vdash
        s=t\rightarrow
        \bigl((\psi[x/s]\square\chi[x/s])
        \leftrightarrow
        (\psi[x/t]\square\chi[x/t])\bigr).
    \]
    
    Resta tratar o caso em que $\varphi$ é $Qy\,\psi$, em que $Q$ é $\forall$ ou $\exists$.

    Se $y=x$, então as substituições por $x$ não atravessam o quantificador inicial.
    Assim,
    \[
        \varphi[x/s]=\varphi=\varphi[x/t],
    \]
    e a conclusão segue de
    \[
        T\vdash s=t\rightarrow(\varphi\leftrightarrow\varphi).
    \]

    Suponha, então, que $y$ é diferente de $x$.
    Por hipótese, $x$ é substituível por $s$ e por $t$ em $Qy\,\psi$, e, portanto, em $\psi$.

    Se $x$ não ocorre livre em $\psi$, então
    \[
        \psi[x/s]=\psi=\psi[x/t],
    \]
    e a conclusão segue imediatamente.

    Suponha, então, que $x$ ocorre livre em $\psi$.
    Como $x$ é substituível por $s$ e por $t$ em $Qy\,\psi$, a variável $y$ não ocorre em $s$ nem em $t$.
    Portanto, $y$ não ocorre livre em $s=t$.

    Pela hipótese de indução,
    \[
        T\vdash s=t\rightarrow(\psi[x/s]\leftrightarrow\psi[x/t]).
    \]
    Aplicando o Metalema~\ref{mlem:quantificacaoEquivalenciasSobHipotese}, com $\theta$ sendo $s=t$, obtemos
    \[
        T\vdash
        s=t\rightarrow
        \bigl(Qy\,\psi[x/s]\leftrightarrow Qy\,\psi[x/t]\bigr).
    \]
\end{proof}

Em linguagem informal, a Lei da igualdade de Leibniz diz que objetos iguais satisfazem as mesmas fórmulas, desde que as substituições envolvidas sejam válidas.

Como regra de inferência derivada, obtemos a forma usual de substituição de iguais por iguais.

\begin{metacorollary}[Substituição de iguais por iguais]\label{mcor:substituicaoIguaisPorIguais}
    Seja $T$ uma teoria de primeira ordem.
    Seja $x$ uma variável, sejam $s$ e $t$ termos da linguagem de $T$, e seja $\varphi$ uma fórmula da linguagem de $T$.
    Suponha que $x$ seja substituível por $s$ e por $t$ em $\varphi$.

    Então vale a seguinte regra de inferência:
    \[
        \frac{T\vdash s=t \qquad T\vdash \varphi[x/s]}
             {T\vdash \varphi[x/t]}.
    \]
\end{metacorollary}

\begin{proof}
    Pelo Metateorema~\ref{mthm:leiIgualdadeLeibniz},
    \[
        T\vdash s=t\rightarrow(\varphi[x/s]\leftrightarrow\varphi[x/t]).
    \]
    Como $T\vdash s=t$, segue que
    \[
        T\vdash \varphi[x/s]\leftrightarrow\varphi[x/t].
    \]
    Em particular,
    \[
        T\vdash \varphi[x/s]\rightarrow\varphi[x/t].
    \]
    Como $T\vdash \varphi[x/s]$, concluímos
    \[
        T\vdash \varphi[x/t].
    \]
\end{proof}

\begin{example}
    Na linguagem da teoria dos conjuntos, suponha que
    \[
        T\vdash a=b.
    \]
    Aplicando a Lei da igualdade de Leibniz à fórmula $z\in u$, com a variável $u$, obtemos
    \[
        T\vdash a=b\rightarrow(z\in a\leftrightarrow z\in b).
    \]
    Portanto,
    \[
        T\vdash z\in a\leftrightarrow z\in b.
    \]
\end{example}

\begin{example}
    Suponha que $a,b,c,u,z$ são variáveis distintas na linguagem da teoria dos conjuntos.
    Aplicando a Lei da igualdade de Leibniz à fórmula
    \[
        \forall z(z\in u\rightarrow z\in c),
    \]
    com a variável $u$, obtemos
    \[
        T\vdash
        a=b\rightarrow
        \Bigl(
        \forall z(z\in a\rightarrow z\in c)
        \leftrightarrow
        \forall z(z\in b\rightarrow z\in c)
        \Bigr).
    \]
    Aqui a hipótese de substituibilidade é satisfeita porque a variável $z$ não ocorre nos termos $a$ e $b$.
\end{example}

\begin{example}
    Na linguagem da teoria dos grupos, suponha que $x,y,z$ são variáveis.
    Pelo Metalema~\ref{mlem:substituicaoIguaisTermos}, aplicado ao termo $x\cdot z$, temos
    \[
        T\vdash x=y\rightarrow x\cdot z=y\cdot z.
    \]
    Analogamente,
    \[
        T\vdash x=y\rightarrow z\cdot x=z\cdot y.
    \]
    Assim, a igualdade é compatível com a multiplicação em cada coordenada.
\end{example}

\subsection*{Exercícios}

\begin{exercise}
    Use a Metaproposição~\ref{mprop:propriedadesElementaresIgualdade} para justificar as seguintes regras derivadas:
    \[
        \frac{T\vdash r=s}{T\vdash s=r}
        \qquad\text{e}\qquad
        \frac{T\vdash r=s \quad T\vdash s=t}{T\vdash r=t}.
    \]
\end{exercise}

\begin{exercise}
    Seja $T$ uma teoria de primeira ordem e sejam $r,s,t$ termos da linguagem de $T$.
    Escolha uma variável $u$ que não ocorre em $r$, $s$ nem $t$.
    Use a Lei da igualdade de Leibniz, aplicada à fórmula $r=u$, para mostrar que
    \[
        T\vdash s=t\rightarrow(r=s\leftrightarrow r=t).
    \]
\end{exercise}

\begin{exercise}
    Na linguagem da teoria dos conjuntos, use a Lei da igualdade de Leibniz para mostrar que
    \[
        T\vdash a=b\rightarrow
        \bigl(a\in c\leftrightarrow b\in c\bigr)
    \]
    e que
    \[
        T\vdash a=b\rightarrow
        \bigl(c\in a\leftrightarrow c\in b\bigr).
    \]
\end{exercise}

\begin{exercise}
    Na linguagem da teoria dos grupos, use o Metalema~\ref{mlem:substituicaoIguaisTermos} para mostrar que
    \[
        T\vdash x=y\rightarrow x^{-1}=y^{-1}.
    \]
    Depois, mostre que
    \[
        T\vdash x=y\rightarrow (x\cdot z)^{-1}=(y\cdot z)^{-1}.
    \]
\end{exercise}